\chapter{Security, Privacy and Access Control}
\section{Authentication and authorisation boundaries}
Passwords use bcrypt hashing. Access-token validation checks request identity, while persistent refresh families and device sessions support rotation and revocation. Browser refresh requires an appropriate origin/custom header when authenticated by a cookie. A Linux-specific cookie name avoids a same-host, different-port collision with the Windows fallback. These controls follow layered session-management principles rather than relying on a JWT alone~\cite{owasp}.

Every file operation resolves current ownership or inherited access. Search receives only readable indexed IDs, and source mapping checks authority again before returning citations. History retrieval additionally checks actor ownership and retained document scope. A private service token authenticates Express to FastAPI; the backend rejects an arbitrary intelligence URL before sending that token. The browser cannot submit a trusted actor header directly to a publicly exposed intelligence endpoint.
Previously issued presigned download URLs remain usable until their short expiry; application revocation does not retroactively invalidate an already issued S3 signature. New authenticated requests are checked against current authority. This distinction prevents a claim of instantaneous revocation for every transfer path.
Figures~\ref{fig:auth} and~\ref{fig:trust} make the identity and service boundaries explicit.
\diagram{authentication}{Authentication and refresh-session sequence. A token identifies a caller; current resource policy still governs each operation.}{auth}
\diagram{trust}{Trust boundaries for browser input, private intelligence, generation and guarded cloud access. Planned instance-role execution remains unverified on EC2.}{trust}
\section{Threats and mitigations}
\par\begin{table}[htbp]\caption{Threat analysis and residual scope}\small\begin{tabularx}{\linewidth}{p{.25\linewidth}YY}\toprule Threat&Implemented mitigation&Remaining qualification\\\midrule
Foreign document exposure&Current file scope; source re-authorisation; history ownership&Regression cases are finite\\
Refresh-token replay&Hashed family records and rotation/revocation&Secure deployment origin still matters\\
Prompt injection&Untrusted evidence delimiters; no model tools; quote checks&No proof against every adversarial prompt\\
Internal token leakage&Approved service-host allowlist&Host administrator remains privileged\\
Stale index or revision&Generation checks, current metadata and outbox&No cross-store distributed transaction\\
Cloud overspend&Free-plan/credit guard; idle event routing&Billing delay prevents a hard cost cap\\
Container compromise&Non-root, dropped capabilities, isolated ports/networks&Host enforcement awaits EC2 verification\\\bottomrule\end{tabularx}\end{table}\par
\section{Actual dependency and image findings}
The final image snapshots contain 23 distinct HIGH/CRITICAL CVE identifiers across all four production images. Package-level counts are larger because one CVE may affect several installed packages or appear in multiple images. The table therefore reports findings without conflating them with unique defects. Production npm audit and a Python advisory scan have different coverage from whole-image scanning; the latter additionally inspects OS packages and bundled Go binaries.
\par\begin{table}[htbp]\caption{Trivy image snapshot: package findings, not unique CVEs}\label{tab:security}\small\begin{tabularx}{\linewidth}{YrrY}\toprule Image&Critical&High&Boundary\\\midrule
API&0&52&OS findings; no high Node-package result\\
Intelligence&0&53&OS findings; no high Python-package result\\
CPU Ollama&0&54&53 OS findings plus one bundled Go finding\\
PostgreSQL&1&72&Stable libxml2 finding remains\\\bottomrule\end{tabularx}\end{table}\par
The initial Python image's SQLite/zlib critical findings were addressed by the base-image change. Unused npm, pip/ensurepip and root-only gosu tooling were removed from final runtime images where appropriate; available stable security updates were applied. The remaining PostgreSQL libxml2 CVE-2026-6653 is scanner-critical, while Debian classifies it as a minor issue with no stable fix~\cite{libxml}. The database has no published port and the application exposes no arbitrary SQL/XML interface. This limits exposure but does not fix the vulnerable package.

Other findings include privileged util-linux prerequisites, SQLite FTS table construction and a vendor Ollama x/net HTML-parser issue. Non-root execution, dropped capabilities, restricted network boundaries and the absence of model web-search use reduce specific attack paths. Findings are retained, not suppressed. No production-security certification or universal absence of vulnerabilities is claimed.
\section{Secrets, logs and cloud controls}
The report package contains no credential files, refresh cookies, private databases or recovery backups. AWS diagnostic values are sanitised and account identifiers are omitted where unnecessary. The proposed EC2 role would use instance credentials and IMDSv2 rather than copied personal profiles. A scoped deny prevents unrelated Parameter Store reads inherited from SSM core permissions. Actual metadata firewall, encrypted disk and private SSM controls remain post-approval checks~\cite{imds,ssmsession}.
\FloatBarrier
